\section{A complete metabelian correspondence}
\label{sec:metabelian}

Assume in this section that the polar Lie algebra is metabelian.
Then \(W=L'\) is abelian, \(f|_W\) is linear, and the induced map on
\(U=V/W\) is linear. Both operators satisfy \(T^2+T=\id\).
We give \(U\) and \(W\) the \(\F_4\)-structures in which these operators
are multiplication by \(\omega\).

\begin{lemma}[The operator pencil]\label{lem:operator-pencil}
In any binary-linear splitting of \(V\to U\), write
\begin{equation}\label{eq:preliminary-form}
 f(x,z)=(\omega x,\omega z+L_xz+q(x)).
\end{equation}
The binary-linear pencil \(x\mapsto L_x\) satisfies
\[
 L_x\omega=\omega L_x,\quad L_x^3=0,\quad
 L_{\omega x}=\omega L_x+L_x^2,\quad
 L_xL_y=L_yL_x,\quad L_x^2L_y=0.
\]
\end{lemma}
\begin{proof}
Put \(J=\omega\id_W\) and \(M_x=J+L_x\).
Both \(M_x\) and \(\id+M_x\) are invertible: they are the restrictions
on coset differences of the permutations \(f\) and \(f+\id=f^2\).
Substitution in \eqref{eq:functional} gives
\[
 M_{\omega x}M_x=\id+M_x.
\]
For \(A=M_x\), the three matrices along the \(\omega\)-orbit are
\(A,\id+A^{-1},(\id+A)^{-1}\). Their sum is \(J\), because the
three \(L\)-terms sum to zero. Thus \(J\) is a rational function of
\(A\), and in particular commutes with it. Clearing the denominators gives
\[
 JA(\id+A)=A^3+A+\id.
\]
Since \(J^2+J=\id\) and \(J^3=\id\), this is
\((A+J)^3=0\). It follows that \(L_x\) commutes with \(J\) and is
cube zero. Expanding \((J+L_x)^{-1}\) now gives
\(L_{\omega x}=JL_x+L_x^2\).
The additivity of the left side forces the \(L_x\)'s to commute.

For \(P=L_x,Q=L_y\), the identity \((P+Q)^3=0\) yields
\(P^2Q=PQ^2=:T\). Apply the same identity to
\(P'=JP+P^2=L_{\omega x}\) and \(Q\). It gives
\((J^2+J)T=P^2Q^2\), hence \(T=P^2Q^2=PT\).
As \(P\) is nilpotent, this implies \(T=0\).
\end{proof}

\begin{theorem}[Normal form and equivalence]\label{thm:correspondence}
Every metabelian quadratic solution is binary-linearly conjugate to
\begin{equation}\label{eq:normal-form}
 f(x,z)=\bigl(\omega x,\omega z+P(x)z+Q(x)+P(x)Q(x)\bigr)
       \quad\text{on }\F_4^d\oplus W,
\end{equation}
where
\begin{align}
 P(x)&=\sum_i x_iP_i,& Q(x)&=\sum_{i<j}x_ix_ja_{ij},\label{eq:PQ}\\
 P_iP_j&=P_jP_i,&P_i^2&=0,\label{eq:square-zero}\\
 P_i a_{jk}+P_j a_{ik}+P_k a_{ij}&=0\quad(i<j<k),
 \label{eq:triangle}\\
 W&=\sum_iP_iW+\Span_{\F_4}\{a_{ij}:i<j\}.
 \label{eq:intrinsic-generation}
\end{align}
Conversely these conditions define a quadratic solution with intrinsic
polar image \(W\).

The binary-linear conjugacy groupoid of these maps is equivalent to the
isomorphism groupoid of finite-dimensional metabelian \(\F_4\)-Lie
algebras \(\mathfrak l\) satisfying
\begin{equation}\label{eq:Lie-action-square}
 (\ad u)^2|_{\mathfrak l'}=0\qquad(u\in\mathfrak l).
\end{equation}
For corresponding objects,
\(C_{\GL(V,2)}(f)\cong\Aut_{\F_4}(\mathfrak l)\).
Their lower central dimensions agree after multiplication by two.
\end{theorem}

\begin{proof}
Lemma~\ref{lem:operator-pencil} provides the preliminary actions.
We first normalize a section and separate the scalar weights of the
cocycle (Steps 1--2). The decisive step is to combine binary degree two
with intrinsic generation to force square zero (Step 3).
We then verify sufficiency, construct the auxiliary Lie algebra and
compare lower central series (Steps 4--5). Finally, tracking changes of
coordinates upgrades the normal form to an equivalence on morphisms
(Step 6).

\emph{Step 1: a normalized section.}
Choose an \(\F_4\)-basis \(e_i\) of \(U\), with arbitrary lifts \(v_i\)
in \(V\). Define a binary-linear section by
\(e_i\mapsto v_i\), \(\omega e_i\mapsto f(v_i)\).
The images form a basis modulo \(W\).
Since \(f(v_i+f(v_i))=v_i\), the cocycle \(q\) in
\eqref{eq:preliminary-form} vanishes on each coordinate \(\F_4\)-line.

\emph{Step 2: scalar weights.}
Initially put \(P_i=L_{e_i}\), without yet asserting \(P_i^2=0\).
Lemma~\ref{lem:operator-pencil} gives
\[
 L_{\alpha e_i}=\alpha P_i+(\alpha+\alpha^2)P_i^2
       \quad(\alpha\in\F_4).
\]
Set \(N(x)=\sum_i x_i(P_i+P_i^2)\).
Then \(L_x=N+N^2\), \(N^3=0\), and
\(L_{\omega x}=\omega N+\omega^2N^2\).
Decompose \(q=q_0+q_1+q_2\) into the three scalar weights
\(q_j(\omega x)=\omega^jq_j(x)\). Projection to each weight is a linear
combination of \(q(x),q(\omega x),q(\omega^2x)\), so preserves binary
degree and the vanishing on coordinate lines.

The cocycle equation is
\[
 q(\omega x)=\omega^2q(x)+(\omega N+\omega^2N^2)q(x).
\]
Its weight components give
\[
 q_0=Nq_2+\omega N^2q_1,\qquad
 q_1=\omega Nq_0+\omega^2N^2q_2.
\]
Using \(N^3=0\) yields
\begin{equation}\label{eq:weight-solution}
 q_0=Nq_2,\qquad q_1=N^2q_2.
\end{equation}

\emph{Step 3: binary degree forces square zero.}
The reduced \(\F_4\)-polynomial of a weight-two binary-quadratic map has
the form \(\sum_i b_ix_i^2+\sum_{i<j}a_{ij}x_ix_j\).
Line normalization removes the first sum, so \(q_2=Q\).
Consequently
\[
 q_1=N^2Q=\sum_{k,i<j}x_k^2x_ix_jP_k^2a_{ij}.
\]
This is where the degree bound becomes decisive. For \(k=i\) or \(j\),
the exponent pattern is \((3,1)\); otherwise it is \((2,1,1)\).
Both patterns have binary degree three, whereas the weight projection
\(q_1\) has degree at most two. The cubed or squared coordinate
identifies \(k\), and the remaining coordinates identify \(\{i,j\}\).
Thus each displayed term has its own reduced monomial, forcing
\begin{equation}\label{eq:squares-kill-coefficients}
 P_k^2a_{ij}=0\qquad\text{for all }k,i,j.
\end{equation}
For completeness, the one-coordinate function basis
\(1,x,x^2,x^3\) has binary degrees \(0,1,1,2\).
Its leading classes form a basis in the associated graded degree
filtration. Tensoring over the coordinates proves that distinct reduced
monomials of the displayed degree three cannot cancel to degree two.
This justifies the comparison in the binary degree filtration.

We now extend this vanishing from the cocycle coefficients to the whole
intrinsic derived ideal. Let \(\m W=\sum_iP_iW\). By the preliminary identities,
every \(P_k^2\) annihilates \(\m W\).
Modulo \(\m W\), the polar values in
\eqref{eq:preliminary-form} are exactly the values of the polar of \(Q\).
Since \(W\) is the span of all polar values, this proves
\eqref{eq:intrinsic-generation}.
Together with \eqref{eq:squares-kill-coefficients}, it gives \(P_k^2=0\)
on all of \(W\). Now \(L_x=P(x)\) and \(q=Q+PQ\).
The degree-three terms of \(PQ\) with three distinct indices vanish
precisely when \eqref{eq:triangle} holds. The other terms have exponent
patterns \((2,1)\), of binary degree two. This proves necessity.

\emph{Step 4: sufficiency.}
Conditions \eqref{eq:square-zero} and \eqref{eq:triangle} make
\(P(x)^2=0\) and make \(PQ\) binary-quadratic.
Under \(x\mapsto\omega x\), the functions \(Q\) and \(PQ\) have weights
two and zero. Thus
\[
 (Q+PQ)(\omega x)=\omega^2Q+PQ
   =\omega^2(Q+PQ)+\omega P(Q+PQ).
\]
Together with
\((\omega\id+\omega P)(\omega\id+P)=\omega^2\id+P\),
this verifies \(f^2=f+\id\) by substitution.
Mixed polar values generate \(\m W\); modulo that space, the values
of the polar of \(Q\) generate the \(a_{ij}\)'s over \(\F_4\).
Equation~\eqref{eq:intrinsic-generation} therefore makes the polar
image exactly \(W\).

\emph{Step 5: the associated Lie algebra.}
Put \(\beta(x,y)=Q(x+y)+Q(x)+Q(y)\) and define
\begin{equation}\label{eq:auxiliary-bracket}
 [(x,z),(y,t)]_0=(0,P(x)t+P(y)z+\beta(x,y)).
\end{equation}
Jacobi follows from commuting operators and \eqref{eq:triangle}.
This is a metabelian \(\F_4\)-Lie algebra with derived ideal \(W\),
and it satisfies \eqref{eq:Lie-action-square}.
Conversely, splitting the abelianization of any algebra with these
properties gives exactly such operators and an alternating cocycle.
Choose its cross-term refinement \(Q\) to recover
\eqref{eq:normal-form}.

For both the original polar bracket and \eqref{eq:auxiliary-bracket},
the action on \(W\) is \(P(x)\). Hence, with \(\m\) the ideal generated
by these actions,
\begin{equation}\label{eq:lower-central-module}
 \gamma_2=W,\qquad\gamma_k=\m^{k-2}W\quad(k\ge3).
\end{equation}
This proves the lower central comparison.

\emph{Step 6: morphisms.}
A binary-linear conjugacy preserves the intrinsic space \(W\).
Its diagonal blocks \(A:U\to U'\) and \(C:W\to W'\) commute with
the intrinsic scalar actions, hence are \(\F_4\)-linear.
Write its off-diagonal block as \(t_s+t_a\), where
\(t_s(\omega x)=\omega t_s(x)\) and
\(t_a(\omega x)=\omega^2t_a(x)\).
Direct substitution in \(Sf=f'S\) gives the two conditions
\begin{align}
 CP(x)&=P'(Ax)C,\label{eq:morphism-action}\\
 Q'(Ax)&=CQ(x)+t_a(x)+P'(Ax)t_s(x).
 \label{eq:morphism-cocycle}
\end{align}
Indeed the unseparated cocycle equation reduces to the second one by
multiplication with \(\id+P'(Ax)\), whose inverse is itself.
Polarization over \(\F_4\) eliminates \(t_a\) and gives
\begin{equation}\label{eq:morphism-Lie}
 \beta'(Ax,Ay)=C\beta(x,y)
              +P'(Ax)t_s(y)+P'(Ay)t_s(x).
\end{equation}
Equations~\eqref{eq:morphism-action} and \eqref{eq:morphism-Lie}
say exactly that \((x,z)\mapsto(Ax,Cz+t_s(x))\) is an
\(\F_4\)-Lie isomorphism.

Conversely, for any such Lie isomorphism the function
\(Q'(Ax)+CQ(x)+P'(Ax)t_s(x)\) has zero polar and is an ordinary
homogeneous quadratic. It is therefore a unique sum of squares
\(\sum b_ix_i^2\). Taking this as \(t_a\) gives a unique conjugacy.
Scalar parts respect composition because \(A\) and \(C\) are
\(\F_4\)-linear. Thus the correspondence is an equivalence on both
objects and morphisms, proving also the centralizer assertion.
\end{proof}

\begin{corollary}\label{cor:metabelian-nilpotent}
Every metabelian polar algebra is nilpotent and satisfies
\((\ad v)^3=0\) for every \(v\). If \(w=\dim_{\F_4}W\), its class is
at most \(w+1\).
\end{corollary}
\begin{proof}
The commuting square-zero operators generate a nilpotent associative ideal.
They have a common nonzero kernel whenever \(W\ne0\); repeating on the
quotient gives a strictly triangular flag, so all products of \(w\)
operators vanish. Apply \eqref{eq:lower-central-module}.
Finally \(\ad v\) maps \(V\) into \(W\), and its square on \(W\) is zero.
\end{proof}

\begin{theorem}\label{thm:solvable-nilpotent}
A polar Lie algebra is solvable if and only if it is nilpotent.
Its solvable radical is invariant under \(f\) and is nilpotent.
\end{theorem}
\begin{proof}
Induct on derived length. The abelian case is immediate.
Put \(W=L'\), \(W'=L''\). The restricted solution on \(W\) has
nilpotent polar algebra by induction. The quotient \(L/W'\) is
metabelian, hence nilpotent by Corollary~\ref{cor:metabelian-nilpotent}.
For fixed \(x\in L\), the derivation \(D=\ad x|_W\) is therefore
nilpotent on \(W/W'\).

We use the following elementary derivation fact. If \(N\) is nilpotent
and \(D\) is nilpotent on \(N/N'\), then it is nilpotent on \(N\).
Each factor \(\gamma_i(N)/\gamma_{i+1}(N)\) is a quotient of
\((N/N')^{\otimes i}\). The induced derivation is the sum of the \(i\)
commuting operators acting in individual tensor positions.
If \(D^s=0\) on \(N/N'\), the sum has power
\(i(s-1)+1\) equal to zero, by the pigeonhole principle on its expanded
products. The finite lower central filtration proves the assertion on
\(N\), in every characteristic.

Apply this fact to \(N=W\). Since \(\ad x\) maps \(L\) into \(W\),
it is nilpotent on \(L\). Engel's theorem now gives nilpotency of \(L\).
The converse is immediate. If \(R\) is the solvable radical,
\(L/R\) has zero center. Lemma~\ref{lem:ideal-descent} makes \(R\)
invariant, and the result just proved applies to the restricted solution.
\end{proof}
